% Clase 19 --- De la espira material a la curva imaginaria (§3.1).
% "Imagínense que ese cable lo hago muy finito, muy finito, muy finito... pero
%  si no hay un circuito, esta fórmula vale igual."
% El paso conceptual de la clase: la ley no habla de cables, habla de curvas.
% Los dos paneles tienen la misma curva y el mismo campo variable; lo único que
% cambia es que en (b) no hay materia.
\begin{center}
\begin{tikzpicture}[scale=1]

  % =============== (a) con cable ===============
  \begin{scope}
    \fill[figblue!8] (0,0) circle (1.9);
    \draw[figgray, line width=0.7pt] (0,0) circle (1.9);
    \foreach \ang in {30,90,150,210,270,330} {
      \node[figblue, scale=0.7] at (\ang:1.1) {$\otimes$};
    }
    \node[figblue, scale=0.7] at (0,0) {$\otimes$};
    \node[figlblaux, anchor=south] at (0,1.95) {$\vec B(t)$ variable};

    \draw[figline, line width=2.2pt] (0,0) circle (1.35);
    \draw[figvecres, line width=1.1pt] (100:1.35) arc (100:65:1.35);
    \node[figlbl, figred, anchor=south west] at (60:1.4) {$I$};
    \node[figlbl, anchor=north, align=center] at (0,-2.75)
      {(a) con un cable: circula corriente};
  \end{scope}

  % =============== (b) sin cable ===============
  \begin{scope}[xshift=6.8cm]
    \fill[figblue!8] (0,0) circle (1.9);
    \draw[figgray, line width=0.7pt] (0,0) circle (1.9);
    \foreach \ang in {30,90,150,210,270,330} {
      \node[figblue, scale=0.7] at (\ang:1.1) {$\otimes$};
    }
    \node[figblue, scale=0.7] at (0,0) {$\otimes$};
    \node[figlblaux, anchor=south] at (0,1.95) {$\vec B(t)$ variable};

    \draw[figamber, dashed, line width=1.3pt] (0,0) circle (1.35);
    \draw[figvecaux, figamber, line width=1pt] (100:1.35) arc (100:65:1.35);
    \node[figlbl, figamber, anchor=south west] at (60:1.4) {$C$};
    \node[figlblaux, anchor=north, align=center] at (0,-1.42)
      {no hay materia:\\[-0.25em] sólo una curva imaginaria};
    \node[figlbl, anchor=north, align=center] at (0,-2.75)
      {(b) sin cable: el campo $\vec E$ está igual};
  \end{scope}

\end{tikzpicture}\\[0.5em]
{\small\itshape Vale la pena detenerse en este paso porque es el que cambia el
estatus de la ley. Mientras haya un cable uno puede contarse el cuento de que
``la corriente inducida circula por el circuito''; pero si el cable se va
haciendo cada vez más fino, y finalmente se lo saca, la ley
$\oint_C \vec E\cdot d\vec l = -d\Phi_B/dt$ sigue valiendo palabra por palabra:
sólo necesita una curva y una superficie apoyada en ella, no materia. Lo que hay
en el espacio, entonces, no es ``corriente esperando un cable'' sino un
\textbf{campo eléctrico inducido}, que existe con o sin circuito y que es el que
efectivamente empuja las cargas cuando el cable está ---el campo magnético no
puede hacerlo, porque la fuerza magnética nunca trabaja---. La consecuencia
teórica es fuerte: este $\vec E$ tiene circulación no nula, así que \emph{no}
deriva de un potencial y $\vec E = -\nabla V$ deja de valer. Siguen valiendo,
eso sí, $\vec F = q\vec E$, el trabajo y la ley de Ohm.}
\end{center}
